\hypertarget{appendix-ea2dadd21ce5}{}
\section{Transformer attention and KV cache runtime}\label{appendix-ea2dadd21ce5-transformer-attention-and-kv-cache-runtime}

A transformer runtime becomes understandable when tensor shapes remain visible. For one attention head, let a token representation produce query, key and value vectors

\[
q,k,v \in \mathbb{R}^{d}.
\]

For a query attending to keys \(k_j\), the scaled score is

\[
\ell_j = \frac{q^T k_j}{\sqrt{d}},
\]

followed by a stable softmax

\[
p_j = \frac{e^{\ell_j-m}}{\sum_t e^{\ell_t-m}},
\qquad m=\max_t \ell_t,
\]

and the output is

\[
o = \sum_j p_j v_j.
\]

Subtracting \(m\) does not change the distribution because the same factor appears in numerator and denominator, but it prevents overflow when logits are large.

\subsection{Causal masking}\label{appendix-ea2dadd21ce5-causal-masking}

For autoregressive decoding at token position \(t\), a query must not observe keys and values from future positions. In a batch attention matrix this is represented by a causal mask before softmax. During incremental decoding, a correctly managed KV cache gives the same structural guarantee: only keys and values through the current token have been appended.

The cache exists because recomputing all earlier K/V projections every token is wasteful. If each layer has key/value vectors of head width \(d_h\), number of KV heads \(h_{kv}\), sequence length \(L\), and element size \(b\), the rough storage per layer is

\[
2 L h_{kv} d_h b
\]

bytes before allocator and alignment overhead. Cache datatype, paging strategy and maximum context length are therefore memory-system decisions.

\subsection{Multi-head and grouped-query attention}\label{appendix-ea2dadd21ce5-multi-head-and-grouped-query-attention}

Multiple heads perform the same mechanism with separate projected subspaces. Grouped-query attention may use fewer key/value heads than query heads; the mapping from query head to KV head must be explicit. GGUF metadata commonly carries head-count and key/value dimension information, so model loading and runtime shape validation meet at this boundary.

\subsection{Residual, normalization, and feed-forward structure}\label{appendix-ea2dadd21ce5-residual-normalization-and-feed-forward-structure}

Attention is only one sublayer. A complete transformer block also needs residual paths, normalization, and a feed-forward network or architecture-specific alternative. The order---pre-norm versus post-norm, RMSNorm versus LayerNorm, activation and gating choice---is part of the model architecture and must come from the model contract rather than a generic ``transformer'' assumption.

\subsection{Executable witness}\label{appendix-ea2dadd21ce5-executable-witness}

tools/labs/phase29\_ml\_runtime\_lab.cpp implements stable softmax and a small incremental KV cache. Two keys with distinct values make the weighting observable; a single-token case proves that incremental causal attention cannot see a future token that has not been appended.

It intentionally does not claim multi-head batching, rotary position embeddings, grouped-query attention, tokenizer behavior, paged KV caches, fused kernels or production-model parity. Those are separate implementation milestones.

Related: \hyperref[appendix-cebadf3106d1]{Tensor runtime layout and views}, \hyperref[appendix-fef72dd246f2]{GGUF binary model loading}, \hyperref[appendix-a87217e3af2b]{ONNX graph import and lowering}, \hyperref[appendix-241c0e42ac2d]{Quantized matvec and dequantization}.
